\newpage
\begin{center}
\large\textbf{Минобрнауки России}

\large\textbf{Юго-Западный государственный университет}
\vskip 1em
\normalsize{Кафедра программной инженерии}
\vskip 1em
\ifВКР{
        \begin{flushright}
        \begin{tabular}{p{.4\textwidth}}
        \centrow УТВЕРЖДАЮ: \\
        \centrow Заведующий кафедрой \\
        \hrulefill \\
        \setarstrut{\footnotesize}
        \centrow\footnotesize{(подпись, инициалы, фамилия)}\\
        \restorearstrut
        «\underline{\hspace{1cm}}»
        \underline{\hspace{3cm}}
        20\underline{\hspace{1cm}} г.\\
        \end{tabular}
        \end{flushright}
        }\fi
\end{center}
\vspace{1em}
  \begin{center}
  \large
\ifВКР{
ЗАДАНИЕ НА ВЫПУСКНУЮ КВАЛИФИКАЦИОННУЮ РАБОТУ
  ПО ПРОГРАММЕ БАКАЛАВРИАТА}
  \else
ЗАДАНИЕ НА КУРСОВУЮ РАБОТУ (ПРОЕКТ)
\fi
\normalsize
  \end{center}
\vspace{1em}
{\parindent0pt
  Студента \АвторРод, шифр\ \Шифр, группа \Группа
  
1. Тема «\Тема\ \ТемаВтораяСтрока»
\ifВКР{
	утверждена приказом ректора ЮЗГУ от \ДатаПриказа\ № \НомерПриказа
}\fi.

2. Срок предоставления работы к защите \СрокПредоставления

3. Исходные данные для создания программной системы:

3.1. Перечень решаемых задач:}

\renewcommand\labelenumi{\theenumi)}

\begin{enumerate}
\item Провести анализ деятельности редакционно-издательского отдела вуза и выявить проблематику традиционного подхода к проверке учебно-методических материалов;
\item Разработать концептуальную модель информационно-аналитической системы;
\item Спроектировать архитектуру системы, включая базу данных и веб-интерфейс;
\item Реализовать систему средствами веб-технологий (Python, FastAPI, PostgreSQL, HTML, CSS, JavaScript);
\item Разработать модуль автоматической проверки оформления документов формата DOCX (шрифт, отступы, интервалы, нумерация страниц);
\item Разработать модуль генерации титульных листов для методических указаний, учебных пособий и монографий.
\end{enumerate}

{\parindent0pt

3.2. Входные данные и требуемые результаты для программы:}
\begin{enumerate}
\item Входными данными для программной системы являются: файлы формата DOCX (методические указания, учебные пособия, монографии), параметры для генерации титульных листов (название, автор, рецензенты, УДК, ISBN, описание), учётные данные пользователей (ФИО, email, пароль, факультет, кафедра, роль).

\item Выходными данными являются: результаты автоматической проверки документа (список ошибок с указанием страниц), исправленный документ с визуально выделенными ошибками, PDF-отчёт о проверке, сгенерированный титульный лист в формате DOCX, статистические отчёты по загруженным документам, JWT-токены для авторизации, подтверждение отправки файлов на электронную почту РИО.
\end{enumerate}

{\parindent0pt

4. Содержание работы (по разделам):

4.1. Введение.

4.2. Анализ предметной области.

4.3. Техническое задание: основание для разработки, назначение разработки,
требования к программной системе, требования к оформлению документации.

4.4. Технический проект: общие сведения о программной системе, проект
данных программной системы (ER-модель, реляционная модель), проектирование архитектуры программной системы, проектирование пользовательского интерфейса программной системы.

4.5. Рабочий проект: маршруты приложения, спецификация контроллеров, системное тестирование разработанного веб-сайта.

4.6. Заключение.

4.7. Список использованных источников.

5. Перечень графического материала:

\списокПлакатов

\vskip 2em
\begin{tabular}{p{6.8cm}C{3.8cm}C{4.8cm}}
	Руководитель \ifВКР{ВКР}\else работы (проекта) \fi & \lhrulefill{\fill} & \fillcenter\Руководитель\\
	\setarstrut{\footnotesize}
	& \footnotesize{(подпись, дата)} & \footnotesize{(инициалы, фамилия)}\\
	\restorearstrut
	Задание принял к исполнению & \lhrulefill{\fill} & \fillcenter\Автор\\
	\setarstrut{\footnotesize}
	& \footnotesize{(подпись, дата)} & \footnotesize{(инициалы, фамилия)}\\
	\restorearstrut
\end{tabular}
}

\renewcommand\labelenumi{\theenumi.}
